\documentclass[11pt,a4paper]{article}

\usepackage[italian]{babel}
\usepackage{booktabs}
\usepackage{graphicx}
\usepackage{geometry}
\usepackage{longtable}
\usepackage{xcolor}
\usepackage{hyperref}

\geometry{margin=2.5cm}
\newcommand{\pending}{\textcolor{gray}{DA COMPLETARE}}
\newcommand{\safefigure}[2]{%
  \IfFileExists{#1}{\includegraphics[width=#2\linewidth]{#1}}{%
    \fbox{\parbox[c][4cm][c]{0.9\linewidth}{\centering
      Figura non ancora generata\\\texttt{\detokenize{#1}}}}}}

\title{Selezione dei frame e riconoscimento dei numeri di maglia\\
Risultati sperimentali}
\author{Football Tracking --- artifact quantitativo per la tesi}
\date{Aggiornato al 21 luglio 2026}

\begin{document}
\maketitle

\section{Protocollo sperimentale}

Tutti i confronti OCR usano le stesse tracklet, lo stesso numero massimo di
crop e lo stesso backend di aggregazione. Gli split sono disgiunti per
sequenza. Le nove sequenze di development possono essere usate per diagnosi e
ablation; le undici sequenze frozen non possono essere usate per modificare
soglie o scegliere checkpoint. Lo split GSR test rimane congelato.

Le metriche aggregate devono essere accompagnate da statistiche per sequenza.
Gli intervalli di confidenza sono calcolati con bootstrap al livello della
sequenza, non del crop, per non assumere indipendenza tra osservazioni della
stessa partita o tracklet.

\section{Dataset e split}

\begin{table}[htbp]
  \centering
  \caption{Split usati negli esperimenti jersey. Una sequenza puo' comparire
  soltanto nel ruolo dichiarato per lo stesso esperimento.}
  \label{tab:dataset-splits}
  \resizebox{\textwidth}{!}{%
  \begin{tabular}{lrrrrl}
    \toprule
    Split & Sequenze & Tracklet & Crop/immagini & Annotati & Uso \\
    \midrule
    GSR train OCR/readability & 37 & 311 & 624 & 624 & training \\
    Development interno & 9 & 75 & 154 & 154 & model selection/ablation \\
    Frozen interno & 11 & 173 & 11072 frame manifest & 0 & evaluation \\
    GSR test & 49 & 774 & 2954 crop selezionati & 0 & final locked evaluation \\
    \bottomrule
  \end{tabular}%
  }
\end{table}

\begin{table}[htbp]
  \centering
  \caption{Derivazione SJN-210k usata per il pretraining specifico. Lo split
  ufficiale test e' impiegato come validation e selezione del checkpoint, non
  come test finale indipendente.}
  \label{tab:sjn-dataset}
  \begin{tabular}{lrrrr}
    \toprule
    Partizione & Immagini & Regioni localizzate & Una cifra & Due cifre \\
    \midrule
    SJN train & 14956 & 8938 & 4782 & 10174 \\
    SJN validation & 6454 & 3789 & 2144 & 4310 \\
    \bottomrule
  \end{tabular}
\end{table}

\section{Test ufficiale GSR}

La configurazione congelata dopo development e frozen e' stata valutata una
sola volta sulle 49 sequenze GSR test: 774 tracklet e 2954 crop selezionati.
Baseline e candidato condividono esattamente selezione dei frame, detector,
1803 regioni rilevate, padding e aggregazione.

\begin{table}[htbp]
  \centering
  \caption{Risultati sul test ufficiale GSR.}
  \label{tab:ctc-test-aggregate}
  \IfFileExists{tables/test_ocr_aggregate.tex}{%
    \resizebox{\textwidth}{!}{\input{tables/test_ocr_aggregate.tex}}%
  }{\pending: generazione da artifact GSR test}
\end{table}

Il CTC SJN$\rightarrow$GSR aumenta le track corrette da 249 a 442 e riduce le
errate da 356 a 163, a copertura invariata. L'accuracy all-track passa da
0.322 a 0.571 ($+0.249$; bootstrap paired sequence-level CI 95\%
$[0.222,0.278]$). La disponibilita' del GT nei top-5 passa da 0.499 a 0.673
($+0.174$; CI 95\% $[0.148,0.202]$).

\begin{table}[htbp]
  \centering
  \caption{Transizioni paired sul test ufficiale. Le 219 track discordanti
  comprendono 206 recuperi e 13 regressioni.}
  \label{tab:ctc-test-transitions}
  \IfFileExists{tables/test_ocr_transitions.tex}{%
    \resizebox{\textwidth}{!}{\input{tables/test_ocr_transitions.tex}}%
  }{\pending: generazione da artifact GSR test}
\end{table}

Il test binomiale esatto a due code sulle track discordanti produce
$p=7.54\times10^{-46}$. Il bootstrap usa 10000 campioni di 49 sequenze intere
con seed 20260720; nessun crop o track viene campionato come osservazione
indipendente.

\begin{table}[htbp]
  \centering
  \caption{Annotazioni manuali per target di leggibilita' del numero.}
  \label{tab:manual-labels}
  \begin{tabular}{lrrr}
    \toprule
    Partizione & Readable & Unreadable & Uncertain/escluse \\
    \midrule
    Totale review & 583 & 195 & 22 \\
    Train & 453 & 171 & 0 \\
    Validation interna & 130 & 24 & 0 \\
    \bottomrule
  \end{tabular}
\end{table}

\paragraph{Distribuzione jersey.}
La tabella completa per numero di maglia verra' generata dal manifest del
benchmark finale e inserita come \texttt{tables/jersey\_distribution.tex}.

\section{Classificatori per la selezione dei frame}

\begin{table}[htbp]
  \centering
  \caption{Risultati interni dei classificatori di frame. I valori non
  disponibili nei checkpoint upstream sono indicati con ``--''.}
  \label{tab:frame-classifiers}
  \resizebox{\textwidth}{!}{%
  \begin{tabular}{llrrrrrr}
    \toprule
    Target & Backbone & Train & Val & Accuracy & Precision & Recall & AP \\
    \midrule
    Legibility legacy & ResNet34 & -- & -- & -- & -- & -- & -- \\
    Clean back & YOLO11s-cls & 484 & 114 & 0.895 & -- & -- & 0.932 \\
    Number readable & YOLO26s-cls & 624 & 154 & 0.981 & 0.992 & 0.985 & 0.999 \\
    \bottomrule
  \end{tabular}%
  }
\end{table}

\begin{table}[htbp]
  \centering
  \caption{Qualita' del ranking sul validation manuale sequence-disjoint.}
  \label{tab:frame-ranking}
  \begin{tabular}{lrrrr}
    \toprule
    Modello & Tracklet eleggibili & Top-5 hit & Top-5 hit rate & Frozen osservate \\
    \midrule
    YOLO11s clean-back & 38 & 38 & 1.000 & 0 \\
    YOLO26s readability & 61 & 61 & 1.000 & 0 \\
    \bottomrule
  \end{tabular}
\end{table}

\section{OCR track-level: ablation interna}

\begin{table}[htbp]
  \centering
  \caption{Ablation 3$\times$2 sulle nove sequenze di development. ``GT top-5''
  conta le tracklet in cui almeno un crop produce il numero ground truth.}
  \label{tab:ocr-internal}
  \resizebox{\textwidth}{!}{%
  \begin{tabular}{llrrrrrr}
    \toprule
    Selettore & Recognizer & Assigned & Correct & Wrong & Coverage & Acc. all & GT top-5 \\
    \midrule
    ResNet34 & SAR pretrained & 117 & 20 & 97 & 0.755 & 0.129 & 82 \\
    ResNet34 & SAR fine-tuned & 136 & 66 & 70 & 0.877 & 0.426 & 145 \\
    YOLO11s clean-back & SAR pretrained & 138 & 20 & 118 & 0.890 & 0.129 & 84 \\
    YOLO11s clean-back & SAR fine-tuned & 155 & 56 & 99 & 1.000 & 0.361 & 148 \\
    YOLO26s readability & SAR pretrained & 134 & 23 & 111 & 0.865 & 0.148 & 86 \\
    YOLO26s readability & SAR fine-tuned & 155 & 59 & 96 & 1.000 & 0.381 & 147 \\
    \bottomrule
  \end{tabular}%
  }
\end{table}

Il risultato mostra che il maggiore incremento deriva dal recognizer SAR
fine-tuned. I classificatori clean-back/readability aumentano la disponibilita'
del candidato corretto solo marginalmente e peggiorano la decisione aggregata.

\section{OCR track-level: validation frozen}

\begin{table}[htbp]
  \centering
  \caption{Ablation del recognizer sui 92 crop development con regioni
  manuali, mantenendo invariata l'aggregazione track-level.}
  \label{tab:ctc-transfer-development}
  \begin{tabular}{lrrrr}
    \toprule
    Recognizer & Crop corretti & Track corrette & Acc. track & GT top-5 \\
    \midrule
    CTC GSR precedente & 61/92 & 31/46 & 0.674 & 40/46 \\
    CTC SJN zero-shot & 71/92 & 38/46 & 0.826 & 45/46 \\
    CTC SJN$\rightarrow$GSR & \textbf{77/92} & \textbf{40/46} & \textbf{0.870} & \textbf{45/46} \\
    \bottomrule
  \end{tabular}
\end{table}

\begin{table}[htbp]
  \centering
  \caption{Confronto sulle undici sequenze frozen. Selezione dei frame,
  detector GSR smoke, confidenza detector 0.25, padding 0.25 e aggregazione
  sono identici; varia soltanto il checkpoint CTC.}
  \label{tab:ocr-frozen}
  \IfFileExists{tables/frozen_ocr_aggregate.tex}{%
    \resizebox{\textwidth}{!}{\input{tables/frozen_ocr_aggregate.tex}}%
  }{\pending: generazione da artifact frozen}
\end{table}

Il pretraining SJN seguito dal fine-tuning GSR porta le decisioni corrette da
84 a 124 su 173 track, con copertura invariata. Le transizioni paired sono 45
wrong-to-correct e 5 correct-to-wrong ($p=4.21\times10^{-9}$, test binomiale
esatto a due code sulle 50 track discordanti). Il checkpoint resta congelato
in audit fino alla singola valutazione GSR test.

\begin{table}[htbp]
  \centering
  \caption{Transizioni paired sul frozen set.}
  \label{tab:ctc-frozen-transitions}
  \IfFileExists{tables/frozen_ocr_transitions.tex}{%
    \resizebox{\textwidth}{!}{\input{tables/frozen_ocr_transitions.tex}}%
  }{\pending: generazione da artifact frozen}
\end{table}

\section{Risultati per sequenza}

\IfFileExists{tables/test_ocr_per_sequence.tex}{%
  \input{tables/test_ocr_per_sequence.tex}
}{%
\begin{table}[htbp]
  \centering
  \caption{Metriche per video. Questa tabella sara' importata automaticamente
  da \texttt{per\_sequence.csv}.}
  \label{tab:per-sequence}
  \begin{tabular}{llrrrrr}
    \toprule
    Sequenza & Metodo & Tracklet & Correct & Wrong & Coverage & Acc. all \\
    \midrule
    \multicolumn{7}{c}{\pending: generazione automatica dal benchmark multi-video} \\
    \bottomrule
  \end{tabular}
\end{table}
}

\section{Runtime e costo computazionale}

\begin{table}[htbp]
  \centering
  \caption{Costo dei componenti. Tempi e memoria devono essere misurati sulla
  stessa GPU e con gli stessi crop.}
  \label{tab:runtime}
  \resizebox{\textwidth}{!}{%
  \begin{tabular}{lrrrrr}
    \toprule
    Componente & Checkpoint MB & Crop & Batch & Tempo totale s & Peak GPU MB \\
    \midrule
    ResNet34 legibility & \pending & \pending & 32 & \pending & \pending \\
    YOLO11s clean-back & 11 & \pending & 32 & \pending & \pending \\
    YOLO26s readability & 11 & \pending & 32 & \pending & \pending \\
    SAR pretrained & \pending & \pending & 8 & \pending & \pending \\
    SAR fine-tuned & 220 & \pending & 8 & \pending & \pending \\
    \bottomrule
  \end{tabular}%
  }
\end{table}

\section{Analisi statistica}

Per ogni metodo saranno riportati media, deviazione standard, mediana e
intervallo interquartile delle metriche per sequenza. Gli intervalli di
confidenza al 95\% saranno stimati con bootstrap paired delle sequenze. Le
transizioni tra baseline e candidato saranno conteggiate sulle stesse tracklet:
recovered-correct, correct-to-wrong, new-wrong-emission, wrong-to-wrong e
abstention.

\IfFileExists{tables/test_ocr_statistics.tex}{%
  \input{tables/test_ocr_statistics.tex}
}{%
\begin{table}[htbp]
  \centering
  \caption{Riepilogo statistico sequence-level.}
  \label{tab:statistics}
  \begin{tabular}{lrrrrrr}
    \toprule
    Metodo & Media acc. & Std & Mediana & IQR & CI 95\% low & CI 95\% high \\
    \midrule
    \multicolumn{7}{c}{\pending: bootstrap sul benchmark multi-video finale} \\
    \bottomrule
  \end{tabular}
\end{table}
}

\section{Figure}

\begin{figure}[htbp]
  \centering
  \safefigure{figures/ocr_accuracy_coverage.png}{0.85}
  \caption{Accuracy e coverage per metodo.}
  \label{fig:accuracy-coverage}
\end{figure}

\begin{figure}[htbp]
  \centering
  \safefigure{figures/ocr_coverage_accuracy_scatter.png}{0.75}
  \caption{Trade-off tra coverage e accuracy.}
  \label{fig:coverage-accuracy}
\end{figure}

\begin{figure}[htbp]
  \centering
  \safefigure{figures/ocr_per_sequence_boxplot.png}{0.8}
  \caption{Distribuzione dell'accuracy tra sequenze.}
  \label{fig:per-sequence}
\end{figure}

\begin{figure}[htbp]
  \centering
  \safefigure{figures/readability_score_distribution.png}{0.8}
  \caption{Distribuzione degli score readable/unreadable.}
  \label{fig:readability-scores}
\end{figure}

\begin{figure}[htbp]
  \centering
  \safefigure{figures/readability_precision_recall.png}{0.75}
  \caption{Curva precision--recall del classificatore readability.}
  \label{fig:precision-recall}
\end{figure}

\begin{figure}[htbp]
  \centering
  \safefigure{figures/topk_curve.png}{0.75}
  \caption{Probabilita' di trovare il numero corretto al variare di $k$.}
  \label{fig:topk}
\end{figure}

\begin{figure}[htbp]
  \centering
  \safefigure{figures/ocr_transition_matrix.png}{0.7}
  \caption{Matrice delle transizioni baseline--candidato.}
  \label{fig:transitions}
\end{figure}

\begin{figure}[htbp]
  \centering
  \safefigure{figures/jersey_confusion_matrix.png}{0.75}
  \caption{Confusion matrix dei numeri di maglia.}
  \label{fig:jersey-confusion}
\end{figure}

\begin{figure}[htbp]
  \centering
  \safefigure{figures/ocr_prediction_bias.png}{0.8}
  \caption{Distribuzione dei numeri predetti e bias ricorrenti.}
  \label{fig:prediction-bias}
\end{figure}

\begin{figure}[htbp]
  \centering
  \safefigure{figures/pipeline_diagram.png}{0.95}
  \caption{Pipeline sperimentale per selezione frame e OCR.}
  \label{fig:pipeline}
\end{figure}

\begin{figure}[htbp]
  \centering
  \safefigure{figures/qualitative_examples.png}{0.95}
  \caption{Esempi qualitativi: recuperi, regressioni e astensioni.}
  \label{fig:qualitative}
\end{figure}

\section{Stato delle decisioni}

Il selettore rimane ResNet34. La catena migliore e validata sul test ufficiale
usa detector GSR smoke e CTC preaddestrato SJN seguito da fine-tuning GSR.
YOLO clean-back, YOLO readability e SAR fine-tuned non sono promossi. Il
benchmark autorizza la promozione del nuovo recognizer nella configurazione
congelata; l'integrazione runtime deve comunque preservare provenance e
supportare audit prima dell'eventuale apply.

\end{document}
